\documentclass[a4paper,11pt]{article}

\usepackage{revy}
\usepackage[utf8]{inputenc}
\usepackage[T1]{fontenc}
\usepackage[danish]{babel}


\revyname{MatematikRevyen}
\revyyear{2026}
\version{0.1}
\eta{$3$ minutter}
\status{Færdig}

\title{Akademisk kvarter}
\author{Flemming '25, Vibe '25 }
\melody{Sabrina Carpenter: ``15 Minutes''}

\begin{document}
\maketitle

\begin{roles}
\role{S}[Sanger] Studerende
\role{F}[Sanger] Forelæser
\end{roles}

\begin{song}
  \sings{S}[Vers1] Jeg hører uret tikke
  Kan ikke nå at smør’ min rugbrødsmad
  Skema sig’r klokken otte
  Løber ned af gaden, farer afsted
  (Er) der på slaget, takket vær’ en smutvej
  Maven rumler, savner leverpostej
  Men der er intet der kan stresse mig


  \sings{S}[Omkv1] For jeg kan gøre meget, på et kvarter
  Købe i kantinen, lave ugeopgaver
  Tag et kort spil whist, du ved at jeg vipper
  Med, med, du ved at
  Jeg kan gøre meget, på et kvarter
  Ud og børste tænder, sludre med mine venner
  Gå en lille tur, hent kaffe op på fjerde
  Du, du, du ved jeg kan

  \sings{F}[Vers2] Klar til dagens for’læsninger
  Pakket noter og mit eget kridt
  Min kop er fyldt med kaffe
  Lig’ kommet ind i Auditorie 1.
  Gud, fejlen på det slide er min egen 
  Et par steder mangler jeg dollartegn
  Og kan se, at LaTex ik’ vil compile

  \sings{F}[Omkv2] Men jeg kan gøre meget, på et kvarter
  Se min KU-wrapped , lave ugeopgaver
  Ordne mine filer, fikse min’ projekter
  Ja, ja, ja jeg kan
  Jeg kan gøre meget på et kvarter
  Ignorer en mail, måske svare sener'
  Tænd min mikrofon, find de batterier
  Du, du, du ved jeg kan

  \sings{S} Ah
  Yeah, ah
  Yeah, ah

  \sings{S+F}[Bro] Akademisk fritid
  Det gir’ mig ekstra frihed
  Jeg kan starte en Risk strid
  Ja jeg kan, det ik’ vanvid
  Akademisk fritid
  Det gir’ mig ny selvtillid
  Adopter et nyt ged’kid
  Ja jeg kan, det ik’ vanvid

  \sings{S+F}[Outro] Uh-uh, uh-uh, uh
  Uh, ja jeg kan
  Oh, ja jeg kan
  Oh, ja jeg kan, babe


\end{song}

\end{document}

